% =====================================================================
%  Applied Machine Learning — Assignment Report
%  University of Sussex, Spring 2026
% ---------------------------------------------------------------------
%  Compile with:
%     pdflatex report
%     bibtex   report
%     pdflatex report
%     pdflatex report
%
%  Or use latexmk (handles all of the above automatically):
%     latexmk -pdf report.tex
%
%  This template uses natbib + BibTeX so it compiles on any TeX install
%  (Overleaf, MiKTeX, MacTeX, TeX Live, Sussex lab machines).
% ---------------------------------------------------------------------
%  Word limit: 3000 (includes figure captions, excludes references)
%  Suggested per-section budget is noted in a comment above each section.
% =====================================================================

\documentclass[11pt,a4paper]{article}

% --- Encoding & fonts -------------------------------------------------
\usepackage[T1]{fontenc}
\usepackage[utf8]{inputenc}

% --- Page layout ------------------------------------------------------
\usepackage[margin=2.5cm]{geometry}
% microtype is loaded with safe options so it works with the default
% Computer Modern font on minimal TeX installs.
\usepackage[protrusion=true,expansion=false]{microtype}
\usepackage{setspace}
\onehalfspacing

% --- Maths ------------------------------------------------------------
\usepackage{amsmath}
\usepackage{amssymb}
\usepackage{bm}

% --- Graphics, tables, captions --------------------------------------
\usepackage{graphicx}
\graphicspath{{figures/}}
\usepackage{booktabs}
\usepackage{array}
\usepackage{tabularx}
\usepackage{multirow}
\usepackage{caption}
\usepackage{subcaption}
\captionsetup{font=small,labelfont=bf}

% --- Floats & placement ----------------------------------------------
\usepackage{float}
\usepackage{placeins}

% --- Hyperlinks (load late) ------------------------------------------
\usepackage{xcolor}
\usepackage[hidelinks]{hyperref}
\usepackage{cleveref}

% --- TikZ for placeholder figures ------------------------------------
\usepackage{tikz}
\usetikzlibrary{positioning,arrows.meta,shapes.geometric}

% --- Bibliography (natbib + BibTeX, universally supported) -----------
\usepackage[numbers,sort&compress]{natbib}
\bibliographystyle{plainnat} % Ships with natbib; works everywhere.
% Other options: unsrtnat (in citation order), abbrvnat (abbreviated).

% --- Custom commands -------------------------------------------------
% Drafting helpers — REMOVE all \todo before submission.
\newcommand{\todo}[1]{\textcolor{red}{\textbf{[TODO:} #1\textbf{]}}}
% Placeholder image so the document compiles before you have figures.
\newcommand{\placeholder}[3][8cm]{%
  \begin{center}
    \begin{tikzpicture}
      \draw[dashed,gray!60,line width=0.6pt] (0,0) rectangle (#1,#2);
      \node[align=center,gray!80,font=\small\itshape] at (#1/2,#2/2) {Placeholder: #3};
    \end{tikzpicture}
  \end{center}
}
% Once a real figure exists in figures/, replace \placeholder with:
% \includegraphics[width=0.8\linewidth]{filename}

% =====================================================================
%  Title block
% =====================================================================
\title{%
  Applied Machine Learning Assignment\\[0.4em]
  \large Sentiment Analysis and Face Alignment%
}
\author{%
  \todo{Your name}\\
  \small Candidate number: \todo{xxxxxx}\\
  \small University of Sussex --- G6061 Applied Machine Learning, Spring 2026%
}
\date{\todo{Submission date}}

\begin{document}
\maketitle

% =====================================================================
%  Optional short abstract.
%  If you keep this, it counts toward your 3000-word limit.
%  Many students drop it — the marking scheme doesn't require it.
% =====================================================================
% \begin{abstract}
% \noindent \todo{Two-to-three sentence summary of both systems and headline results.}
% \end{abstract}

% =====================================================================
%  Notes for the writer (delete before submission)
% ---------------------------------------------------------------------
%  Word budget: 3000 total, including figure captions, excluding refs.
%  Suggested allocation:
%    Task 1 Methods                ~600 words
%    Task 1 Results & analysis     ~500 words
%    Task 1 NLTK external eval     ~250 words
%    Task 2 Methods                ~600 words
%    Task 2 Results & analysis     ~500 words
%    Task 2 Robustness analysis    ~250 words
%    Total                         ~2700, with buffer for AI declaration
%
%  Run `texcount report.tex` to track word count as you write.
% =====================================================================

% #####################################################################
\section{Task 1: Sentiment Analysis}
% Marks: 15 (methods) + 15 (analysis) + 10 (NLTK) + 10 (test accuracy)
% #####################################################################

% ---------------------------------------------------------------------
\subsection{Description of methods}
% Marks: 15.  Suggested ~600 words.
% Cover: pre-processing, features/representation, models, hyperparameters,
% compute. Diagrams expected for (a) pre-processing/feature pipeline and
% (b) model architecture.
% ---------------------------------------------------------------------

\subsubsection{Problem framing}

The task is binary sentiment classification of movie reviews, complicated
by an unknown amount of email-style spam evenly distributed across both
labelled classes. \todo{State concretely: the label set the system outputs
is \(\{0, 1, d\}\) where \(d=\) your dummy label.}

\subsubsection{Pre-processing}

\todo{Describe and justify each step: tokenisation, case normalisation,
number normalisation, punctuation handling, stopword decisions
(esp.\ negation words), and either stemming~\cite{porter} or
lemmatisation. Cite NLTK~\cite{nltk}.}

\begin{figure}[H]
  \centering
  \placeholder[14cm]{2.8cm}{Pre-processing flowchart: raw text \(\to\) tokens \(\to\) normalised tokens \(\to\) feature vector}
  \caption{Pre-processing pipeline applied to all training, validation
  and test documents. \todo{Replace placeholder with real flowchart and
  rewrite this caption to describe the choices, e.g. ``Negation tokens
  retained; lemmatisation preferred over stemming following the
  trade-off discussed in Section~X.''}}
  \label{fig:t1-preproc}
\end{figure}

\subsubsection{Features and representations}

\todo{Describe the at-least-two representations you compared, e.g.\ TF--IDF
versus averaged Word2Vec embeddings~\cite{mikolov2013word2vec,mikolov2013distributed}.
Explain how each was computed and why.}

The TF--IDF score for term \(t\) in document \(d\) is
\begin{equation}
  \mathrm{tfidf}(d,t) = \mathrm{tf}(d,t)\,\log\!\left(\frac{N}{\mathrm{df}(t)}\right),
  \label{eq:tfidf}
\end{equation}
where \(N\) is the corpus size and \(\mathrm{df}(t)\) is the number of
documents containing \(t\)~\cite{jurafsky2026}.

\subsubsection{Spam handling}

\todo{Describe your spam-filtering stage. The lectures support
cosine-similarity-to-class-centroid filtering and word-list-based
filtering; both are defensible. State the threshold and how you chose it.}

The cosine similarity between two vectors \(\bm{a}\) and \(\bm{b}\) is
\begin{equation}
  \mathrm{sim}(\bm{a},\bm{b}) = \frac{\bm{a}\cdot\bm{b}}{\sqrt{(\bm{a}\cdot\bm{a})(\bm{b}\cdot\bm{b})}},
  \label{eq:cosine}
\end{equation}
which we use to compare a candidate document against the centroids of
the declared positive and negative training classes.

\subsubsection{Sentiment models}

\todo{Describe each model you tried: word-list baseline, Naive
Bayes~\cite{jurafsky2026}, Logistic Regression, and (if used)
RNN~\cite{jurafsky2026}. For each: ML task type, model family, loss
function, hyperparameters and how chosen.}

For a binary Logistic Regression classifier with weight vector
\(\bm{w}\), the probability of the positive class is
\begin{equation}
  P(y{=}1\mid\bm{x}) = \sigma(\bm{w}^\top\bm{x}) = \frac{1}{1+\exp(-\bm{w}^\top\bm{x})},
\end{equation}
and parameters are obtained by minimising the cross-entropy loss.

\begin{figure}[H]
  \centering
  \placeholder[14cm]{4cm}{Model pipeline: pre-processed text \(\to\) features \(\to\) spam filter \(\to\) sentiment classifier}
  \caption{End-to-end sentiment system. \todo{Rewrite caption to describe
  your final system at a high level.}}
  \label{fig:t1-system}
\end{figure}

\subsubsection{Hyperparameters and compute}

\begin{table}[H]
  \centering
  \small
  \caption{Hyperparameters and how chosen for each Task 1 approach.}
  \label{tab:t1-hyperparams}
  \begin{tabular}{lll}
    \toprule
    Approach & Hyperparameter & Chosen value (method) \\
    \midrule
    Word-list baseline & List size, margin \(\delta\) & \todo{} \\
    TF--IDF + Naive Bayes & N-gram range, smoothing \(\alpha\) & \todo{} \\
    TF--IDF + Logistic Regression & Regularisation \(C\) & \todo{} \\
    Spam threshold & Cosine cut-off & \todo{(swept on validation)} \\
    \bottomrule
  \end{tabular}
\end{table}

\todo{One short paragraph on compute: hardware (Colab T4 / your laptop),
training time per model, peak memory if relevant.}

% ---------------------------------------------------------------------
\subsection{Analysis of results}
% Marks: 15.  Suggested ~500 words.
% Quantitative (confusion matrix required) + qualitative (subjective
% interpretation of language). Failure cases and bias.
% ---------------------------------------------------------------------

\subsubsection{Quantitative evaluation}

\begin{figure}[H]
  \centering
  \begin{subfigure}[t]{0.32\linewidth}
    \placeholder[5cm]{4cm}{CM Approach 1}
    \caption{\todo{Approach name}}
  \end{subfigure}\hfill
  \begin{subfigure}[t]{0.32\linewidth}
    \placeholder[5cm]{4cm}{CM Approach 2}
    \caption{\todo{Approach name}}
  \end{subfigure}\hfill
  \begin{subfigure}[t]{0.32\linewidth}
    \placeholder[5cm]{4cm}{CM Approach 3}
    \caption{\todo{Approach name}}
  \end{subfigure}
  \caption{Confusion matrices on the held-out validation set.
  \todo{Write a sentence or two pointing the reader at the diagonal,
  the spam row, and the most common confusion.}}
  \label{fig:t1-confmats}
\end{figure}

\begin{table}[H]
  \centering
  \small
  \caption{Validation-set performance across approaches.
  Precision, recall and F1 are macro-averaged across the three classes.}
  \label{tab:t1-results}
  \begin{tabular}{lcccc}
    \toprule
    Approach & Accuracy & Precision & Recall & F1 \\
    \midrule
    Word-list baseline (no spam handling) & \todo{} & \todo{} & \todo{} & \todo{} \\
    TF--IDF + NB (no spam handling)       & \todo{} & \todo{} & \todo{} & \todo{} \\
    TF--IDF + LR + cosine spam filter     & \todo{} & \todo{} & \todo{} & \todo{} \\
    \todo{Word2Vec + LR + cosine spam}    & \todo{} & \todo{} & \todo{} & \todo{} \\
    \bottomrule
  \end{tabular}
\end{table}

\todo{Discuss the precision--recall trade-off on the spam threshold and
include the curve.}

\subsubsection{Qualitative evaluation and failure cases}

\todo{Pick three to five interesting validation-set errors. For each,
quote a short fragment and explain in one or two sentences \emph{why}
the model failed. Look for: negations, sarcasm, mixed sentiment, genre
vocabulary used descriptively, very short reviews, spam written in a
movie-review style.}

\subsubsection{Systematic bias}

\todo{Identify and discuss systematic biases observed: e.g.\ length bias
on the spam filter, vocabulary bias on the sentiment classifier, class
imbalance effects after spam removal.}

% ---------------------------------------------------------------------
\subsection{Performance on the NLTK external dataset}
% Marks: 10.  Suggested ~250 words.
% ---------------------------------------------------------------------

\todo{Report sentiment accuracy and (importantly) how often the spam
detector fires on \texttt{nltk.corpus.movie\_reviews}, which contains no
spam. Discuss generalisation: vocabulary drift, domain shift, label-noise
differences.}

% =====================================================================
\FloatBarrier
\section{Task 2: Face Alignment}
% Marks: 15 (methods) + 15 (analysis) + 10 (robustness) + 10 (test acc)
% =====================================================================

% ---------------------------------------------------------------------
\subsection{Description of methods}
% Marks: 15.  Suggested ~600 words.
% ---------------------------------------------------------------------

\subsubsection{Problem framing}

The task is to predict the \((x,y)\) coordinates of five facial landmarks
(left eye, right eye, nose, left mouth corner, right mouth corner;
indices 0--4) for each face image. The training data exhibit \todo{briefly
describe the variability you observed: rotations, brightness changes,
noise, etc.}

\subsubsection{Pre-processing}

\todo{Describe each step: resizing (and the matching transformation of
the landmark coordinates), grayscale conversion, intensity scaling,
optional histogram equalisation~\cite{jurafsky2026}.}

\begin{figure}[H]
  \centering
  \placeholder[14cm]{2.8cm}{Pre-processing flowchart for images and matching landmark transforms}
  \caption{Image pre-processing pipeline. Landmark coordinates are
  transformed identically so they remain consistent with the resized
  image geometry.}
  \label{fig:t2-preproc}
\end{figure}

\subsubsection{Features and representations}

\todo{Describe the representations you compared: raw / flattened pixels,
HOG descriptors~\cite{lowe2004sift}, and CNN-learned features.}

A brief reminder of HOG: the image region is broken into overlapping
\(6\times6\) cells, each described by a histogram of gradient orientations
binned over \(0\)--\(180^\circ\); blocks of \(3\times3\) cells are then
concatenated and locally normalised, yielding good invariance to
brightness changes and small rotations.

\subsubsection{Approaches compared}

\textbf{Baseline 0 --- mean face.}
Predict the per-landmark mean of the training set for every test image.
This is the zero-component PCA shape model.

\textbf{Approach 1 --- HOG features + linear regression.}
\todo{Describe.}

\textbf{Approach 2 --- Cascaded regression with pose-indexed features
\cite{cao2014esr,xiong2013sdm}.}
\todo{Describe the iterative refinement: at each stage, recompute HOG
patches at the current landmark estimates and regress the residual.}

\textbf{Approach 3 --- CNN direct coordinate regression.}
\todo{Describe the architecture (conv blocks, kernel size, pooling,
fully-connected head). State that the loss is the squared Euclidean
distance \(\|\bm{h}(\bm{x})-\bm{y}\|^2\), and the optimiser used.}

\todo{Optionally: a fourth approach using heatmap regression and
soft-argmax, with the hourglass-style architecture~\cite{newell2016hourglass}.}

\begin{figure}[H]
  \centering
  \placeholder[14cm]{6cm}{CNN architecture diagram (input \(\to\) conv blocks \(\to\) FC \(\to\) 10 outputs)}
  \caption{Architecture of the convolutional network used for direct
  coordinate regression. \todo{Replace with your own diagram and state
  layer dimensions in the caption or in a small inline table.}}
  \label{fig:t2-cnn}
\end{figure}

\subsubsection{Data augmentation}

\todo{Describe the augmentations applied during training: rotations,
scales, translations, brightness/contrast changes, and horizontal flips
\emph{with a corresponding swap of the left/right landmark indices}
(eyes 0\(\leftrightarrow\)1; mouth corners 3\(\leftrightarrow\)4; nose
unchanged).}

\begin{figure}[H]
  \centering
  \placeholder[14cm]{4cm}{Augmentation pipeline flowchart}
  \caption{Data augmentation pipeline applied online during CNN
  training. The same geometric transform is applied to both the image
  and the landmark coordinates, and on horizontal flips the left/right
  landmark indices are swapped.}
  \label{fig:t2-augment}
\end{figure}

\subsubsection{Hyperparameters and compute}

\begin{table}[H]
  \centering
  \small
  \caption{Hyperparameters and how chosen for each Task 2 approach.}
  \label{tab:t2-hyperparams}
  \begin{tabular}{lll}
    \toprule
    Approach & Hyperparameter & Chosen value (method) \\
    \midrule
    HOG + linear regression  & Cell size, block size & \todo{} \\
    Cascaded regression      & Number of stages \(K\) & \todo{} \\
    CNN direct regression    & LR, batch size, epochs & \todo{} \\
    \bottomrule
  \end{tabular}
\end{table}

\todo{One short paragraph: hardware (Colab T4), per-model training time.}

% ---------------------------------------------------------------------
\subsection{Analysis of results}
% Marks: 15.  Suggested ~500 words.
% Required: cumulative error distribution OR boxplots.
% ---------------------------------------------------------------------

The error metric is the per-landmark Euclidean distance, normalised by
the inter-eye distance \(d_{\text{ie}} = \|\bm{p}_0 - \bm{p}_1\|\):
\begin{equation}
  \mathrm{NME}_i = \frac{\|\hat{\bm{p}}_i - \bm{p}_i\|}{d_{\text{ie}}},
  \label{eq:nme}
\end{equation}
with the per-image error reported as the mean over the five landmarks.

\begin{figure}[H]
  \centering
  \begin{subfigure}[t]{0.48\linewidth}
    \placeholder[7cm]{5cm}{CED curve: all approaches}
    \caption{Cumulative error distribution.}
    \label{fig:t2-ced}
  \end{subfigure}\hfill
  \begin{subfigure}[t]{0.48\linewidth}
    \placeholder[7cm]{5cm}{Boxplot per approach}
    \caption{Per-image NME boxplots.}
    \label{fig:t2-box}
  \end{subfigure}
  \caption{Quantitative comparison on the held-out validation set.
  \todo{Note that the mean face is included as the floor; comment on
  which approach is uniformly better and where the curves cross.}}
  \label{fig:t2-quant}
\end{figure}

\begin{figure}[H]
  \centering
  \placeholder[14cm]{4cm}{Per-landmark error bars across approaches}
  \caption{Per-landmark mean NME for each approach. \todo{Comment on
  which landmark is hardest and why --- e.g.\ the nose has less local
  texture, or mouth corners shift with expression.}}
  \label{fig:t2-perlandmark}
\end{figure}

\subsubsection{Qualitative examples}

\begin{figure}[H]
  \centering
  \begin{subfigure}[t]{0.32\linewidth}
    \placeholder[5cm]{5cm}{Success 1}
    \caption{}
  \end{subfigure}\hfill
  \begin{subfigure}[t]{0.32\linewidth}
    \placeholder[5cm]{5cm}{Success 2}
    \caption{}
  \end{subfigure}\hfill
  \begin{subfigure}[t]{0.32\linewidth}
    \placeholder[5cm]{5cm}{Success 3}
    \caption{}
  \end{subfigure}\\[0.5em]
  \begin{subfigure}[t]{0.32\linewidth}
    \placeholder[5cm]{5cm}{Failure 1}
    \caption{}
  \end{subfigure}\hfill
  \begin{subfigure}[t]{0.32\linewidth}
    \placeholder[5cm]{5cm}{Failure 2}
    \caption{}
  \end{subfigure}\hfill
  \begin{subfigure}[t]{0.32\linewidth}
    \placeholder[5cm]{5cm}{Failure 3}
    \caption{}
  \end{subfigure}
  \caption{Predicted (red) and ground-truth (green) landmarks on
  representative successes (top) and failures (bottom). \todo{In the
  caption, briefly describe what makes each failure hard.}}
  \label{fig:t2-examples}
\end{figure}

% ---------------------------------------------------------------------
\subsection{Extended robustness analysis}
% Marks: 10.  Suggested ~250 words.
% ---------------------------------------------------------------------

\todo{Describe the perturbations applied (e.g.\ Gaussian noise at
\(\sigma\in\{0,5,10,20,40\}\) on a 0--255 scale; rotations of
\(\{0,5,10,20,40\}^\circ\)) and which approach was tested.}

\begin{figure}[H]
  \centering
  \begin{subfigure}[t]{0.48\linewidth}
    \placeholder[7cm]{5cm}{CED at increasing noise levels}
    \caption{Robustness to additive Gaussian noise.}
  \end{subfigure}\hfill
  \begin{subfigure}[t]{0.48\linewidth}
    \placeholder[7cm]{5cm}{Augmented vs non-augmented training}
    \caption{Effect of training-time augmentation.}
  \end{subfigure}
  \caption{Robustness analysis. \todo{Discuss why the chosen approach is
  more robust to certain perturbations: HOG's invariance to brightness
  and small rotations; convolutional translation equivariance;
  augmentation teaching the network seen transformations.}}
  \label{fig:t2-robust}
\end{figure}

% =====================================================================
\section*{Generative AI declaration}
% Required by the brief.
% =====================================================================

\todo{State concretely how generative AI was used (e.g.\ for ideation,
debugging, or LaTeX formatting), which tool, and that you verified all
output. Per the brief, generative AI was \emph{not} used to write this
report or to analyse data end-to-end.}

% =====================================================================
%  References — excluded from the 3000-word count.
% =====================================================================
\bibliography{references}

\end{document}
